\clearpage
\fancyhead[L]{\small KOMET.}
\hypertarget{native-472-lower}{}
\begin{center}
\rule{0.55\linewidth}{0.4pt}\\[2pt]
Astronomische Nachrichten. Band XXII. Nr. 514. S. 165. 166. 18. Aug. 1844.\\[2pt]
\rule{0.55\linewidth}{0.4pt}
\end{center}
\begin{center}\textsc{Gauss} an \textsc{Schumacher}. Göttingen, 18. August 1844.\end{center}

Ich sende Ihnen noch ein paar Kometenbeobachtungen. Bei dem durchweg ungünstigen Wetter habe ich selbst diese dürftigen Beobachtungen nur mit Mühe machen können:

\begin{center}
\setlength{\tabcolsep}{4pt}
\begin{tabular}{lccc}
1844. & \shortstack{Mittlere Zeit\\in Göttingen} & Rektaszension & \shortstack{Nördliche\\Deklination} \\
Aug. 11 & $9^{\rm u}\ 28^{\rm m}\ 16^{\rm s}$ & $211^\circ\ 34'\ 13''6$ & ---\quad ---\quad --- \\
 & $9\ 31\ 48$ & ---\quad ---\quad --- & $21^\circ\ 48'\ 12''6$ \\
17 & $8\ 34\ 23$ & $208\ 42\ 2.3$ & ---\quad ---\quad --- \\
 & $8\ 36\ 8$ & ---\quad ---\quad --- & $17\ 28\ 22.6$ \\
\end{tabular}
\end{center}

\Needspace{6\baselineskip}
\SourceLeaf{473}
Am 11. dienten drei Sterne, die dem Kometen folgten, zum Vergleich. Der zweite wird mit dem Stern identisch sein, der in \textsc{Bessels} Zone 460 wie folgt angegeben ist
\[
\begin{array}{c|c|c|c}
8.\,9 & 4 & 14^{\rm u}\ 7^{\rm m}\ 50^{\rm s}54 & 21^\circ\ 50'\ 21''9
\end{array}
\]
\noindent sofern man annimmt, dass die Beobachtung am dritten statt am vierten Faden gemacht wurde und die Uhrzeit deshalb irrtümlich um das Fadenintervall vermindert abgedruckt ist. Meine oben angegebene Rektaszension beruht auf dieser Hypothese, deren Richtigkeit ich noch nicht bestätigen konnte, und auf der Annahme, dass die von \textsc{Bessel} verwendete Reduktionszahl $15^{\rm s}46$ oder die unmittelbar beobachtete Antrittszeit $14^{\rm u}\ 8^{\rm m}\ 6^{\rm s}0$ war. Würde man den zweiten Stern ganz ausschließen, wäre die Rektaszension des Kometen um $4''2$ größer und müsste folglich mit $211^\circ\ 34'\ 17''8$ angesetzt werden. Die Parabel von Dr. \textsc{Goldschmidt} stimmt noch immer recht gut mit den Beobachtungen überein.

Am 17. August wurde Piazzi XIII. 280 zum Vergleich verwendet. Die Rektaszension beruht jedoch nur auf einem einzigen Vergleich.

